\definecolor{azure}{RGB}{52,177,201} % Define the orange color used for highlighting throughout the document


%----------------------------------------------------------------------------------------
%	MAKETITLE DEFINITION 
%----------------------------------------------------------------------------------------

%\newcommand*{\logoleft}[1]{\gdef\@logoleft{#1}}
%\newcommand*{\logorigth}[1]{\gdef\@logorigth{#1}}

\makeatletter
\def\maketitle{
  \begin {titlepage}
    \let \footnotesize \small 
    \let \footnoterule \relax 
    \let \footnote \thanks 
    \@bsmarkseries 
  
%    \begin{figure}[!htb]
%      \begin{minipage}{0.48\textwidth}
%        \centering
%        \includegraphics[width=.8\linewidth]{\@logoleft}
%      \end{minipage}\hfill
%      \begin {minipage}{0.48\textwidth}
%        \centering
%        \includegraphics[width=.9\linewidth]{\@logorigth}
%      \end{minipage}
%    \end{figure}   

    \def \@makefnmark {\rlap {\@textsuperscript {\normalfont \@bsthanksheadpre \tamark \@bsthanksheadpost }}}\long
    \def \@makefntext {\makethanksmark} 
    %\null \vfil  \vskip 5\p@
%    \maketitlehookb {\@bspreauthor \@notenumber \@bspostauthor }
%    \maketitlehookb {\@bspreauthor Version \@versionnumber \@bspostauthor }

    
%    \null \vfil \vskip 1\p@ \vspace *{\droptitle }
    \maketitlehooka {\@bspretitle \Huge\textbf\@title \@bsposttitle } 
    \null \vfil \vskip 10\p@ \vspace *{\droptitle }
    
    \begin{figure}[!h]
      \centering
      \includegraphics[width=.3\linewidth]{oxfordlogo}
    \end{figure}
    
    \null \vfil \vskip 20\p@ \vspace *{\droptitle }
    \maketitlehookb {\@bspreauthor \huge\@author \@bspostauthor } 
    \maketitlehookb {\@bspreauthor St Hilda's College \@bspostauthor } 
    \maketitlehookb {\@bspreauthor University of Oxford \@bspostauthor } 
%    \vspace{\fill}
    \null \vfil \vskip 51\p@ \vspace *{\droptitle }
    \maketitlehookb {\@bspreauthor A thesis submitted in fulfilment of the requirements for the degree of \@bspostauthor } 
    \maketitlehookb {\@bspreauthor \emph{Doctor of Philosophy} \@bspostauthor } 
    \maketitlehookb {\@bspreauthor Hilary Term, 2019 \@bspostauthor }
%    \maketitlehookc {\@bspredate \@date \@bspostdate } 
%    \maketitlehookd \par \@thanks \vfil \null 
  \end{titlepage}\@bscontmark
}
\makeatother


%----------------------------------------------------------------------------------------
%	TABLE OF CONTENTS/FIGURES/TABLES EDITS
%----------------------------------------------------------------------------------------
\makeatletter

\patchcmd{\tableofcontents}
  {\@mkboth{\MakeUppercase\contentsname}{\MakeUppercase\contentsname}}
  {\addcontentsline{toc}{section}{\contentsname}\markboth{\sffamily\normalsize\bfseries\contentsname}{}}
  {}{}
  
\patchcmd{\listoffigures}
  {\@mkboth{\MakeUppercase\listfigurename}{\MakeUppercase\listfigurename}}
  {\addcontentsline{toc}{section}{\listfigurename}\markboth{\sffamily\normalsize\bfseries\listfigurename}{}}
  {}{}
  
\patchcmd{\listoftables}
  {\@mkboth{\MakeUppercase\listtablename}{\MakeUppercase\listtablename}}
  {\addcontentsline{toc}{section}{\listtablename}\markboth{\sffamily\normalsize\bfseries\listtablename}{}}
  {}{}
  
  
\makeatother


%----------------------------------------------------------------------------------------
%	PAGE HEADERS
%----------------------------------------------------------------------------------------

\usepackage{fancyhdr} % Required for header and footer configuration

\pagestyle{fancy}
\renewcommand{\chaptermark}[1]{\markboth{\sffamily\normalsize\bfseries\chaptername\ \thechapter.\ #1}{}} % Chapter text font settings
\renewcommand{\sectionmark}[1]{\markright{\sffamily\normalsize\thesection\hspace{5pt}#1}{}} % Section text font settings
%\fancyhf{} \fancyhead[L,RO]{\sffamily\normalsize\thepage} % Font setting for the page number in the header
\fancyhead[LO]{\rightmark} % Print the nearest section name on the left side of odd pages
\fancyhead[R]{\leftmark} % Print the current chapter name on the right side of even pages
\renewcommand{\headrulewidth}{0.5pt} % Width of the rule under the header
\addtolength{\headheight}{2.5pt} % Increase the spacing around the header slightly
\renewcommand{\footrulewidth}{0pt} % Removes the rule in the footer
\fancypagestyle{plain}{\fancyhead{}\renewcommand{\headrulewidth}{0pt}} % Style for when a plain pagestyle is specified
%\fancyhf{}\fancyfootoffset[F]{3cm}
%\fancyfoot[L]{\thepage}

% Removes the header from odd empty pages at the end of chapters
\makeatletter
\renewcommand{\cleardoublepage}{
\clearpage\ifodd\c@page\else
\hbox{}
\vspace*{\fill}
\thispagestyle{empty}
\newpage
\fi}


%----------------------------------------------------------------------------------------
%	DEFINITION OF COLOURED BOXES
%----------------------------------------------------------------------------------------

\RequirePackage[framemethod=TikZ]{mdframed} % Required for creating the theorem, definition, exercise and corollary boxes


\mdfdefinestyle{exampledefault}{
  rightline=true,
  %frametitlerule=true,
  %frametitlefont=\ttfamily,
  %frametitlerulecolor=black,
  %frametitlebackgroundcolor=lightgray,
  %frametitlerulewidth=2pt,
  backgroundcolor=light-gray,
  roundcorner=5pt,
  %outerlinecolor=blue!70!black,
  %outerlinewidth=2,
  leftmargin=20,
  rightmargin=20,
  ntheorem=true,
  innerleftmargin=15, 
  innertopmargin=1,
  innerbottommargin=1,
  outerlinewidth=1,
  linecolor=light-gray
}
\mdfsetup{
  frametitlealignment=\flushleft
}


% Theorem box
\newmdenv[skipabove=7pt,
skipbelow=7pt,
backgroundcolor=black!5,
linecolor=azure,
innerleftmargin=5pt,
innerrightmargin=5pt,
innertopmargin=5pt,
leftmargin=0cm,
rightmargin=0cm,
innerbottommargin=5pt]{tBox}

% Exercise box	  
\newmdenv[skipabove=7pt,
skipbelow=7pt,
rightline=false,
leftline=true,
topline=false,
bottomline=false,
backgroundcolor=azure!10,
linecolor=azure,
innerleftmargin=5pt,
innerrightmargin=5pt,
innertopmargin=5pt,
innerbottommargin=5pt,
leftmargin=0cm,
rightmargin=0cm,
linewidth=4pt]{eBox}	

% Definition box
\newmdenv[skipabove=7pt,
skipbelow=7pt,
rightline=false,
leftline=true,
topline=false,
bottomline=false,
linecolor=azure,
innerleftmargin=5pt,
innerrightmargin=5pt,
innertopmargin=0pt,
leftmargin=0cm,
rightmargin=0cm,
linewidth=4pt,
innerbottommargin=0pt]{dBox}	

% Corollary box
\newmdenv[skipabove=7pt,
skipbelow=7pt,
rightline=false,
leftline=true,
topline=false,
bottomline=false,
linecolor=gray,
backgroundcolor=black!5,
innerleftmargin=5pt,
innerrightmargin=5pt,
innertopmargin=5pt,
leftmargin=0cm,
rightmargin=0cm,
linewidth=4pt,
innerbottommargin=5pt]{cBox}

% Creates an environment for each type of theorem and assigns it a theorem text style from the "Theorem Styles" section above and a colored box from above
\newenvironment{theorem}{\begin{tBox}\begin{theoremeT}}{\end{theoremeT}\end{tBox}}
\newenvironment{exercise}{\begin{eBox}\begin{exerciseT}}{\hfill{\color{azure}\tiny\ensuremath{\blacksquare}}\end{exerciseT}\end{eBox}}				  
\newenvironment{definition}{\begin{dBox}\begin{definitionT}}{\end{definitionT}\end{dBox}}	
\newenvironment{example}{\begin{exampleT}}{\hfill{\tiny\ensuremath{\blacksquare}}\end{exampleT}}		
\newenvironment{corollary}{\begin{cBox}\begin{corollaryT}}{\end{corollaryT}\end{cBox}}	

%----------------------------------------------------------------------------------------
%	REMARK ENVIRONMENT
%----------------------------------------------------------------------------------------

\newenvironment{remark}{\par\vspace{10pt}\small % Vertical white space above the remark and smaller font size
\begin{list}{}{
\leftmargin=35pt % Indentation on the left
\rightmargin=25pt}\item\ignorespaces % Indentation on the right
\makebox[-2.5pt]{\begin{tikzpicture}[overlay]
\node[draw=azure!60,line width=1pt,circle,fill=azure!25,font=\sffamily\bfseries,inner sep=2pt,outer sep=0pt] at (-15pt,0pt){\textcolor{azure}{R}};\end{tikzpicture}} % Orange R in a circle
\advance\baselineskip -1pt}{\end{list}\vskip5pt} % Tighter line spacing and white space after remark





%----------------------------------------------------------------------------------------
%	SECTION NUMBERING IN THE MARGIN
%----------------------------------------------------------------------------------------

%\makeatletter
%\renewcommand{\@seccntformat}[1]{\llap{\textcolor{azure}{\csname the#1\endcsname}\hspace{1em}}}

\makeatletter
\renewcommand{\@seccntformat}[1]{\textcolor{mygray2}{\csname the#1\endcsname}\hspace{1em}}


%----------------------------------------------------------------------------------------
%	COLORING LINES IN TABLES
%----------------------------------------------------------------------------------------

\makeatletter

\newcommand*{\@rowstyle}{}

\newcommand*{\rowstyle}[1]{% sets the style of the next row
  \gdef\@rowstyle{#1}%
  \@rowstyle\ignorespaces%
}

\newcolumntype{=}{% resets the row style
  >{\gdef\@rowstyle{}}%
}

\newcolumntype{+}{% adds the current row style to the next column
  >{\@rowstyle}%
}

\makeatother




\makeatletter


\newenvironment{alwayssingle}{%
       \@restonecolfalse
       \if@twocolumn\@restonecoltrue\onecolumn
       \else\if@openright\cleardoublepage\else\clearpage\fi
       \fi}%
       {\if@restonecol\twocolumn
       \else\newpage\thispagestyle{empty}\fi}


\newenvironment{abstractseparate} {%
\newgeometry{margin=3.1cm}
\thispagestyle{empty}
\begin{alwayssingle}
 \begin{center}
    {\Large \bfseries \@title \par}
    {{\large \vspace*{2ex} \@author} \par}
    {\large \vspace*{1ex}
    {{St. Hilda's College} \par}
    {University of Oxford \par}
    \vspace*{1ex}
    {{A thesis submitted in fulfilment of the requirements for the degree of} \par}
    {\it {Doctor of Philosophy} \par}
    \vspace*{2ex}
    {Hilary Term, 2019} \par}
%    \vfill
    {\Large \vspace*{1ex}\bfseries  Abstract}
 \end{center}
\vspace{1ex}
   \setlength{\baselineskip}{18pt}
}
{\vfill\end{alwayssingle}\restoregeometry}






\newenvironment{originality} {%
\begin{alwayssingle}
 \begin{center}
    {\Large \bfseries Statement of Originality \par}
 \end{center}
\vspace{1ex}
   \setlength{\baselineskip}{18pt}
}
{\end{alwayssingle}}



\newenvironment{acknowledgements} {%
\begin{alwayssingle}
 \begin{center}
    {\Large \bfseries Acknowledgements \par}
 \end{center}
\vspace{1ex}
   \setlength{\baselineskip}{18pt}
}
{\end{alwayssingle}}

\makeatother


%\renewcommand{\section}{\@startsection{section}{1}{\z@}
%{-4ex \@plus -1ex \@minus -.4ex}
%{1ex \@plus.2ex }
%\normalfont\Large\bfseries}}
%\renewcommand{\subsection}{\@startsection {subsection}{2}{\z@}
%-3ex \@plus -0.1ex \@minus -.4ex}
%{0.5ex \@plus.2ex }
%{\normalfont\sffamily\bfseries}}
%\renewcommand{\subsubsection}{\@startsection {subsubsection}{3}{\z@}
%{-2ex \@plus -0.1ex \@minus -.2ex}
%{.2ex \@plus.2ex }
%{\normalfont\small\sffamily\bfseries}}                        
%\renewcommand\paragraph{\@startsection{paragraph}{4}{\z@}
%{-2ex \@plus-.2ex \@minus .2ex}
%{.1ex}
%{\normalfont\small\sffamily\bfseries}}

%\usepackage{titlesec} 

%\titleformat{\chapter}[display]
%  {\bfseries\Large}
%  {\filright\MakeUppercase{\chaptertitlename} \Huge\thechapter}
%  {1ex}
%  {\titlerule\vspace{1ex}\filleft}
%  [\vspace{1ex}\titlerule]

%\titleformat{\chapter}[display]
%  {\sffamily\bfseries\Huge}
%  {\filleft\MakeUppercase{\chaptertitlename} \Huge\thechapter}
% {1ex}
% {\vspace{1ex}\filleft}
% [\vspace{1ex}]
 







